\documentclass[conference,10pt]{IEEEtran}
\usepackage{booktabs}
\usepackage{amsmath}
\usepackage{url}
\usepackage{array}
\usepackage{balance}
\Urlmuskip=0mu plus 1mu

\title{MCPTrace: Process-Scoped and Verifiable Recovery of Model Context Protocol Artifacts from Memory Images}
\author{\IEEEauthorblockN{Taewook Kwon}
\IEEEauthorblockA{Gyeongbuk Software Meister High School\\
Uiseong, Republic of Korea}}

\begin{document}
\maketitle

\begin{abstract}
Model Context Protocol (MCP) traffic may remain in volatile memory after an agent invokes a tool, but a recovered JSON object alone does not establish its protocol origin, transmission, or execution. Prior work, particularly MCPRecon, has demonstrated memory-based reconstruction of MCP activity. This paper presents MCPTrace, an open-source Volatility 3 plugin that extracts MCP/JSON-RPC candidates from Windows and Linux process address spaces while preserving the original bytes, process identifier, virtual address, encoding, and SHA-256 digest. Its physical-RAM evaluation here is on Linux x86-64, not Windows. It reports unreadable memory and scan limits and distinguishes MCP candidates from LSP, generic JSON-RPC, and ambiguous objects. Across 27 archived Node and .NET HTTP/SSE client captures, an independent server log recorded 342 transmitted messages; an independent memory probe found exact wire bytes for 251. MCPTrace carved all 251 in this fixture. The published MCPRecon CLI carved 184 from a derived stream of the same readable client bytes. These counts include valid JSON-RPC objects regardless of MCP classifier label and are not general detection rates. In a separate four-image Linux full-RAM experiment, the packaged MCP Inspector host exchanged 18 messages over HTTP and stdio; all nine whose exact bytes remained in client memory were carved. Two same-modality LSP and generic JSON-RPC controls produced zero MCP labels. One present Inspector initialization response was carved without an MCP label, while a constructed generic JSON-RPC counterexample was labeled as MCP. These results support address-verifiable extraction and explicit limits on protocol-origin inference.
\end{abstract}

\begin{IEEEkeywords}
memory forensics, Model Context Protocol, Volatility 3, JSON-RPC, tool validation
\end{IEEEkeywords}

\section{Introduction}
Tool-using AI clients communicate with local or remote servers using the Model Context Protocol (MCP). The protocol uses JSON-RPC 2.0 and supports both stdio and HTTP transports \cite{mcp-spec,jsonrpc}. An investigator may need to recover a tool name, call arguments, or a returned value when ordinary logs are missing. Volatile memory is one source of such artifacts, but the interpretation is difficult: a buffer can contain a transmitted message, an SDK object that was never sent, a stale copy from an earlier run, or a similar-looking object from another protocol.

Satter \emph{et al.} introduced MCPRecon, showing that protocol-aware carving can recover MCP interaction traces from memory across clients and transports \cite{mcprecon-paper}. Their published workflow also retains PID, VMA range, and byte offset. Thus, neither MCP-shaped carving nor process-location metadata alone is new here. The present work asks whether exported bytes can be reread at a stated process virtual address in the acquired image, and whether evaluation can separate scanner misses from messages no longer present as exact wire bytes. MCPTrace is a Volatility 3 plugin for that purpose; it does not claim to reconstruct an authoritative event history from memory alone.

This paper makes three contributions. First, it implements an auditable Windows/Linux process scan that exports original bytes, virtual addresses, hashes, and scan coverage, with direct image rereads validating the Linux results. Second, it defines a conditional evaluation that distinguishes transmission, exact-wire presence in readable client memory, and logical-message extraction. Third, it reports a same-capture comparison with the public MCPRecon CLI, a separate live-host full-RAM experiment, and non-MCP controls, including a counterexample to treating MCP-candidate status as proof of protocol origin. The comparison is intended to expose differences and limitations, not to establish universal superiority.

\section{Related Work and Scope}
MCPRecon's published method uses Volatility 3 to enumerate processes and extract VMAs, then indexes retained artifacts by PID, VMA, and byte offset \cite{mcprecon-paper}. It validates JSON-RPC and MCP semantics and associates requests with responses. Its study includes real clients and both stdio and HTTP. The public CLI tested here emits file offsets, parsed JSON, heuristic confidence, and optional raw text; it can accept process-list and region inputs \cite{mcprecon-code}. MCPTrace therefore cannot claim novelty from MCP carving, transport coverage, or process provenance alone.

MCPTrace runs inside Volatility 3 process address translation \cite{volatility}. For Windows it scans readable Virtual Address Descriptors (VADs), and for Linux it scans readable Virtual Memory Areas (VMAs). A record retains the original bytes and their process virtual address, allowing a direct reread from the same image. It also reports unreadable ranges. The comparison validates each tested output in its own representation; a packed-stream file offset is not an original virtual address. This distinction concerns the tested public CLI output, not a claim that MCPRecon's published method lacks process provenance.

MCP and the Language Server Protocol (LSP) both use JSON-RPC and can share fields such as \texttt{clientInfo} and \texttt{capabilities} in an \texttt{initialize} request. LSP may use \texttt{Content-Length} framing and workspace-specific fields \cite{lsp}. MCP's legacy handshake carries a protocol version; the 2026-07-28 revision introduces \texttt{server/discover} and request metadata \cite{mcp-spec}. A syntactically valid JSON-RPC object is consequently insufficient evidence of MCP origin. MCPTrace labels candidates using protocol-specific and conflicting signals; the labels are analytical hints, not findings that a tool was executed.

\section{MCPTrace Design}
\subsection{Acquisition and process scope}
MCPTrace consumes an acquired image through Volatility 3. The plugin enumerates selected processes and their readable address ranges. It requests memory in bounded chunks and carries overlap across contiguous reads so a JSON object may cross a chunk boundary. It does not combine bytes across an unreadable gap or independent VADs/VMAs. The output summary includes regions scanned, bytes read, unreadable bytes, invalid regions, limit events, and errors. A zero-finding scan with extensive unreadable memory must not be described as a complete negative examination.

The scanner searches for JSON-RPC version markers in UTF-8 and UTF-16LE representations, then attempts bounded JSON object recovery. It validates complete objects before classification. The maximum object size and scan bounds are configurable. The plugin also supports partial fragments, but the evaluation in this paper uses complete logical messages. HTTP \texttt{data:} prefixes and stdio line boundaries are contextual bytes when present in the captured region; the scanner does not require a captured transport frame to recover an object.

\subsection{Classification and evidence export}
The classifier considers MCP-specific methods and metadata, legacy protocol-version values, LSP workspace and framing indicators, and conflicts between these signals. Its labels include \texttt{mcp-candidate}, \texttt{lsp-candidate}, \texttt{generic-json-rpc}, and \texttt{ambiguous}. A response may be a candidate because of its structure or relationship to a request, but request/response linkage is conservative: an ID match alone, especially across unrelated processes or sessions, does not prove an interaction. MCPTrace exports the exact carved bytes, their SHA-256 digest, length, encoding, PID, virtual address, and parsed object. The original bytes, rather than a reserialized JSON object, are the target of byte verification.

The evidentiary distinction is important. Let $T$ be the set of transmitted logical messages in the controlled fixture, as recorded by a server or proxy outside the scanned process. Let $P\subseteq T$ be messages whose exact wire bytes were independently located in readable client VMAs at acquisition time, in UTF-8 or UTF-16LE. Let $C$ be the distinct canonical JSON objects emitted by a carver in that scope, regardless of its MCP classifier label. Intersection uses logical JSON equality, not equality of the probe's and carver's byte locations. The conditional coverage reported here is
\begin{equation}
R_{\mathrm{present}}=\frac{|P\cap C|}{|P|}.
\label{eq:present}
\end{equation}
This quantity is not field recall: a transmitted message may be absent from readable memory, retained in another serialization, or outside the examined process. Likewise, $C$ alone does not prove membership in $T$. The external log is available only for the controlled experiments; an ordinary forensic case would not have this ground truth.

\section{Experimental Method}
\subsection{Archived physical-memory cases}
The primary comparison reuses 27 Linux x86-64 QEMU physical RAM captures created with a fresh client process for each run. One condition uses Node 22.14.0 and the official TypeScript MCP SDK 1.32.1 over HTTP with SSE responses. Two conditions use .NET 10.0.404 and the official C\# MCP client SDK 2.3.1, with legacy 2025-11-25 and modern 2026-07-28 protocol flows. The .NET peer is a purpose-built local HTTP fixture, not an official C\# server. Each condition was captured three times at each nominal delay of 0, 5, and 30 seconds after readiness detection. The VM was reused, so memory outside the current client can contain stale data.

A server outside the client logged each actual HTTP JSON body or SSE payload. The original experiment saved a SHA-256 digest for every image, traffic log, and MCPTrace export. It independently probed the same readable client VMAs for each exact wire payload and reread all 626 exported MCPTrace artifact records at their reported virtual addresses. One compressed base RAM image and 26 lossless patches were retained. The pinned MCPTrace source was rerun on each reconstructed, hash-verified image; its exports were compared byte-for-byte with the saved outputs, and new records were reread at their process virtual addresses.

\subsection{MCPRecon configurations and fair scope}
We pinned the public MCPRecon CLI to commit \texttt{4811d04} and verified its full source Git blob. It used the default keywords and 16 KiB window; \texttt{--emit-raw} was enabled solely to validate its output bytes. MCPRecon's default output includes valid JSON-RPC objects without an MCP label; MCPTrace commit \texttt{ded4c00} used \texttt{--include-rpc} for the same counting scope. Two MCPRecon analyses were kept separate. In its native mode, it scanned each whole ELF RAM image. In the scope-matched mode, Volatility reread the current client PID's readable VMA chunks; contiguous chunks were joined and 64 KiB of NUL bytes separated gaps and VMA boundaries. The reread byte, unreadable byte, and VMA counts matched the original independent probe for every image. A second pass checked all 342 traffic rows: each UTF-8 and UTF-16LE exact-wire presence flag in the derived stream matched the original probe. MCPRecon then scanned that stream. Its \texttt{--regions-file} option tests raw file offsets, so virtual VMA addresses cannot be passed directly as ranges in the QEMU ELF image. Packing makes client-byte scope explicit, but changes layout, file offsets, and possibly session clustering. It is not MCPRecon's native whole-image workflow.

For both tools, logical messages were compared with the separately logged traffic using canonical JSON equality, counting repeated copies once per transmitted message. A truth match found anywhere in whole RAM was not automatically attributed to the client. We did not use whole-image unmatched records as a false-positive count, because other processes and previous VM activity can contain legitimate MCP data.

\subsection{Non-MCP controls}
Four previously preserved process-memory captures provide non-MCP controls: LSP and ordinary JSON-RPC clients on macOS ARM64 and Linux ARM64. The macOS captures contain sampled resident private writable pages from an isolated process; the Linux captures were taken from a stopped process. They are process-memory snapshots, not physical RAM images. Their hashes and independent traffic logs were checked. MCPTrace was rerun over their recorded chunks, and MCPRecon was run over a derived stream of the same chunks, with separators at noncontiguous regions.

We also ran both scanners over the exact byte buffers of an existing 16-case constructed protocol corpus. Eight cases have explicitly non-MCP protocol origin, including LSP framing and workspace fields, shared initialization fields, generic JSON-RPC, conflicting signals, and a generic response containing a \texttt{tools} list. These small cases test classification boundaries; their counts cannot estimate a deployment false-positive probability.

\subsection{Live host and controls in physical RAM}
To join live-host and full-RAM validation, we ran the packaged MCP Inspector \texttt{mcpdo} 2.15.0 \cite{inspector} in a fresh Ubuntu 24.04 x86-64 QEMU VM booted with kernel 6.8.0-146. One daemon connected by Streamable HTTP and another by stdio to an official TypeScript-SDK server. The server or an out-of-process proxy journaled actual JSON bodies; each host listed tools and called one. Separate LSP and generic JSON-RPC TCP clients supplied non-MCP controls, with the LSP client using \texttt{Content-Length} framing. Each of the four clients remained alive during one QMP whole-RAM capture. The VM was reused, but only the current client PID was scanned with \texttt{--include-rpc}. The same independent exact-wire probe and virtual-address reread method were applied to each capture. The four compressed images were hash checked and retained privately; call-to-acquisition delay was not synchronized or measured.

\section{Results}
\subsection{Conditional message recovery}
Table~\ref{tab:positive} gives the scope-matched client-memory result. Of 342 logged transmissions, 251 exact wire payloads were located in readable client VMAs. All 251 had UTF-8 copies; 108 also had UTF-16LE copies, and none were UTF-16LE-only. MCPTrace emitted all 251 logical messages in that conditional set; its pinned rerun reproduced all saved exports exactly and reread 626 records at their virtual addresses. MCPRecon's default CLI emitted 184 from the derived stream. Its 315 records were all verified at their packed-stream file offsets, and every record matched a logged logical message. These are carving counts across all classifier labels: MCPTrace labeled 233 of its 251 distinct matches as MCP candidates, whereas MCPRecon labeled 94 of its 184 matches as MCP. The label rules differ, so the classification counts are not a common detection-rate comparison.

\begin{table}[t]
\caption{Same-client-byte comparison on Linux physical RAM captures. ``Present'' means exact wire bytes in readable client VMAs.}
\label{tab:positive}
\centering
\footnotesize
\setlength{\tabcolsep}{3pt}
\begin{tabular}{lrrrrr}
\toprule
Client & Capt. & Sent & Present & Trace & Recon \\
\midrule
Node HTTP/SSE & 9 & 117 & 112 & 112 & 49 \\
.NET legacy & 9 & 117 & 71 & 71 & 68 \\
.NET modern & 9 & 108 & 68 & 68 & 67 \\
\midrule
Total & 27 & 342 & 251 & 251 & 184 \\
\bottomrule
\end{tabular}
\end{table}

The whole-image MCPRecon run matched 252 logged messages. Its 184 matches within $P$ were exactly the same logical messages as the derived-stream matches in every capture; the other 68 lacked exact wire bytes in the current client's readable VMAs. They may be server copies, earlier activity retained by the reused VM, or other representations; this experiment does not assign their source. Thus the whole-image 252 and client-scoped 251/184 figures answer different questions and must not be placed in one ranking.

By requested delay after readiness detection, the 0, 5, and 30 s cohorts each logged 114 messages. Exact wire bytes remained for 86, 86, and 79, respectively; MCPTrace carved those 86, 86, and 79, while MCPRecon carved 63, 64, and 57. The readiness poll was not synchronized to the final call, so these small groups do not establish a decay curve.

\subsection{Observed misses and protocol controls}
Of the 67 exact-wire-present messages not matched by MCPRecon in the derived client stream, 63 occurred in Node cases: nine \texttt{initialize} requests and 54 responses. The other four were three .NET \texttt{notifications/initialized} messages and one \texttt{server/discover} message. On one Node capture, enlarging MCPRecon's window to 64 KiB and 1 MiB left its six-of-thirteen match count unchanged. Inspection of the exact transmitted bytes and the public CLI's \texttt{carve\_json()} shows a concrete mechanism: the algorithm chooses the nearest opening brace before a \texttt{jsonrpc} anchor, which can select a nested object when that key follows other objects. This observation is tied to the tested serialization and CLI configuration; it is not a claim that all MCPRecon configurations fail in this manner.

Table~\ref{tab:negative} separates captured controls from constructed boundary cases. Both tools recovered all eight logged non-MCP messages in the four process snapshots and emitted zero MCP labels there. In the constructed corpus, both labeled one of the eight non-MCP-origin cases as MCP: a generic JSON-RPC response containing a \texttt{tools} list. Thus zero observed labels in the four captured controls cannot be restated as zero false positives in general.

\begin{table}[t]
\caption{Non-MCP controls. Counts are records for captured controls and cases for constructed buffers.}
\label{tab:negative}
\centering
\footnotesize
\setlength{\tabcolsep}{3pt}
\begin{tabular}{lrrr}
\toprule
Control & Input & Trace MCP & Recon MCP \\
\midrule
4 process snapshots & 43 records & 0 & 0 \\
8 constructed cases & 8 cases & 1 & 1 \\
\bottomrule
\end{tabular}
\end{table}

The constructed UTF-16LE LSP buffer was carved by MCPTrace and labeled as LSP; MCPRecon's default CLI did not carve it. This is an encoding boundary case, not evidence of a population-level difference. The unrelated \texttt{tools} response illustrates the opposite boundary: a candidate can be correctly extracted while its protocol origin is incorrectly inferred if the label is treated as definitive.

\subsection{Independent host observations}
The 27 primary captures use controlled SDK clients. A separate earlier experiment ran the packaged MCP Inspector \texttt{mcpdo} host on macOS and Linux ARM64 with stdio and HTTP connections. These are live process-memory captures, not the Linux x86-64 physical RAM specimens used for Table~\ref{tab:positive}. Table~\ref{tab:inspector} reports their original results to make that distinction explicit. The server or a separate stdio proxy logged actual exchanges outside the client. Across four fresh host processes, 36 messages were logged, 22 exact wire payloads remained in the sampled process memory, and all 22 were carved. The remainder were not proven scanner misses; they were not found as complete wire bytes in the sampled regions. No MCPRecon run on these four positive host captures is included in the primary head-to-head comparison.

\begin{table}[t]
\caption{Separate MCP Inspector process-memory observations. These are not physical RAM comparisons.}
\label{tab:inspector}
\centering
\footnotesize
\setlength{\tabcolsep}{4pt}
\begin{tabular}{lrrr}
\toprule
Host / transport & Sent & Present & Carved \\
\midrule
macOS / HTTP & 9 & 3 & 3 \\
macOS / stdio & 9 & 7 & 7 \\
Linux / HTTP & 9 & 4 & 4 \\
Linux / stdio & 9 & 8 & 8 \\
\midrule
Total & 36 & 22 & 22 \\
\bottomrule
\end{tabular}
\end{table}

The four non-MCP controls in Table~\ref{tab:negative} were acquired in the same earlier process-memory campaign, but from separate LSP and JSON-RPC clients. They test those specific settings, not the Linux x86-64 physical-RAM modality. The following experiment supplies same-modality controls, although still only one capture of each.

\subsection{Live host in Linux physical RAM}
Table~\ref{tab:live-ram} reports four new full-RAM captures. The two independent Inspector daemons exchanged 18 messages in total. Exact wire bodies remained in their readable VMAs for nine; MCPTrace carved all nine and reread all 16 emitted artifact records at their virtual addresses. The two non-MCP controls retained all four logged bodies; all four were carved, while none of their 12 verified records carried an MCP label. No scan limit was hit. Of the nine present Inspector messages, eight had at least one MCP-labeled copy. The HTTP \texttt{initialize} response was carved but labeled generic JSON-RPC. Thus the new result validates an end-to-end live host and same-modality control workflow while exposing one classification miss; it does not extend the 251-versus-184 MCPRecon comparison, which was not rerun on these images.

\begin{table}[t]
\caption{Separate Linux physical-RAM captures of live MCP Inspector and non-MCP clients. ``MCP'' counts logged messages with an MCP-labeled copy, not all carved records.}
\label{tab:live-ram}
\centering
\footnotesize
\setlength{\tabcolsep}{3pt}
\begin{tabular}{lrrrr}
\toprule
Current client & Sent & Present & Carved & MCP \\
\midrule
Inspector / HTTP & 9 & 4 & 4 & 3 \\
Inspector / stdio & 9 & 5 & 5 & 5 \\
LSP / TCP & 2 & 2 & 2 & 0 \\
JSON-RPC / TCP & 2 & 2 & 2 & 0 \\
\bottomrule
\end{tabular}
\end{table}

These four images contain 1536 MiB of guest RAM each and are distinct from the archived 27-image SDK comparison and earlier ARM64 process snapshots. The server is purpose-built despite using the official SDK. The VM was reused and each condition was captured once, so the table cannot estimate recovery or false-positive rates across applications or time delays.

\subsection{Why the default CLI missed some present messages}
The Node pattern offers a concrete, testable explanation rather than only a total count. The separately logged JSON texts often place \texttt{jsonrpc} after a nested \texttt{params} or \texttt{result} object. The pinned MCPRecon CLI searches for that marker and chooses the nearest preceding opening brace as a potential JSON start. For the representative Node capture, applying that routine to the isolated transmitted bytes selected an inner object rather than the outer JSON-RPC message for the \texttt{initialize} request and all six responses. These objects were among the seven truth messages absent from its six-of-thirteen client-stream result. Increasing the window did not repair the nearest-brace choice in that capture.

This diagnosis has two boundaries. First, the raw process can hold other serializations of the same logical message; MCPRecon recovered some Node \texttt{tools/call} messages despite that wire-order issue. Second, the pinned public CLI is one implementation and configuration of the MCPRecon work, not a full reimplementation of every procedure in its published paper. The failure analysis supports a narrow claim about this code path and fixture. A broader parser study would need systematically varied key order, nesting depth, escaping, and partial memory objects, with each tool's parser separately tuned.

\section{Discussion and Limitations}
The principal result is a reproducible image-level byte audit with explicit scan coverage, not process provenance by itself. MCPRecon's paper already retains PID, VMA, and offset metadata. MCPTrace's exported bytes, PID, virtual address, and digest let a reader test a record against the acquired image. The 27 SDK captures yielded 626 such reread records; the separate live-host/control RAM experiment yielded 28. This study also verified 315 MCPRecon CLI records at offsets in the derived client streams. A packed offset can be mapped back to a virtual address with the packing index retained by this study, but that original address is not carried in the CLI's raw-file record.

\subsection{Worked image-to-wire audit}
In the public \texttt{node-d30-r2} capture, an external journal logged 13 JSON bodies. An independent probe found eight exact wire bodies in readable client VMAs, and MCPTrace carved all eight. The \texttt{tools/call} request with ID 2 was in the journal, but neither its UTF-8 wire body nor its UTF-16LE form was found in the scanned VMAs. Its SSE response was carved at PID 4361, virtual address \texttt{0x2a99cd19}; 91 bytes were reread from the original image, hash checked, and matched to the external journal. MCPTrace marked this response \texttt{unpaired}. The scan reported 77,991,936 readable and 11,742,543,872 unreadable VMA bytes, with no limit hit; these are virtual-range counts, not physical-RAM size. Thus the response is byte-verifiable, while the absent exact-wire request is not an established parser miss or evidence that no request was sent. The case metadata and checks are public \cite{mcptrace-case}.

\subsection{Interpretation of the denominators}
Three different denominators answer three different questions. The 342 logged messages describe traffic observed by the external server. The 251 exact-wire-present messages describe what the independent memory probe could find in readable client VMAs at acquisition. The 626 exported MCPTrace records count multiple memory copies, not distinct transmitted messages. The 315 MCPRecon records from packed streams likewise count copies. Counting records as messages would inflate coverage; dividing by all 342 transmissions would mix memory retention with parser behavior. Equation~\ref{eq:present} isolates parser recovery conditional on exact-wire presence, but remains sensitive to the probe's encoding and region choices.

The 91 transmissions outside $P$ were not matched by either client-scoped output even under canonical JSON equality, but they are not established false negatives for either scanner. A client's heap may retain partial or reordered objects even when the exact sent bytes are gone. Conversely, a message found in memory can appear in a log buffer, a response cache, or an unsent object. The independent server log allows a logical match in this fixture; it does not identify which specific byte copy was actually transmitted. This is why the output retains candidate status and why the paper does not label every unpaired or unmatched artifact as malicious activity.

\subsection{Reproducibility and data boundaries}
The source, scripts, numeric manifests, and image hashes are documented in the MCPTrace repository \cite{mcptrace}. The 27 comparison RAM archives, byte-identical synthetic external traffic logs, and client-byte replay procedure are public in the research branch. A fresh serial replay of all 27 archives reproduced the 251-versus-184 comparison; per-image numeric reports are provided. The four new live-host/control RAM archives and journals remain private because guest memory may contain ephemeral keys or daemon state; their image-level truth matches cannot be publicly recomputed.

Every reported MCPRecon output record from the comparison was reread at its offset in the exact stream it scanned. The same check does not prove that a record belongs to the client in the whole-image run. The client-only stream maps to readable VMAs, but its offsets are derived and do not preserve the raw ELF file layout. The new MCPTrace exports reproduced the originals byte-for-byte, and their virtual addresses were reread through Volatility's translation from each reconstructed image. Both verification paths depend on the correctness of image acquisition and memory translation.

Several constraints limit inference. The 27 comparison captures come from one reused VM and a purpose-built HTTP server; delays are measured from readiness detection, not the exact last call. The 251-message denominator includes only exact-wire-present content in readable client memory; the other 91 transmissions may have had different in-memory serializations. The added Inspector and same-modality controls are one capture per condition in a separate reused VM, not a random host or field sample. Their call-to-acquisition delay was not measured, every scan reported unreadable VMA bytes, and one present MCP response lacked an MCP label. The earlier four non-MCP process snapshots remain methodologically separate. There is no same-image Windows MCPRecon comparison or estimate of false-positive frequency across software populations. The public MCPRecon CLI may be tuned or modified; this comparison pins one published implementation and its default settings. Packed client streams alter the layout used by its session heuristic. MCPTrace's classifier and scan limits may also change in later versions.

The controls show why investigators should report candidates and uncertainty rather than infer behavior from a JSON object. A recovered \texttt{tools/call} argument is evidence that those bytes existed in memory. It is not, without independent context, proof that a request was sent, accepted, or executed. Request IDs can be reused, and memory copies can outlive the event that created them. The controlled traffic logs make such distinctions measurable here, but would often be unavailable in an incident.

\balance
\section{Conclusion}
MCPTrace implements process-scoped, address-verifiable extraction of MCP/JSON-RPC candidates for Windows and Linux images. This paper's physical-RAM validation is on Linux. In a controlled 27-capture HTTP/SSE fixture, the pinned MCPTrace rerun carved all 251 transmitted messages whose exact wire bytes remained in readable client memory; the pinned public MCPRecon CLI carved 184 from a derived stream of those same bytes. Both figures count valid objects regardless of MCP label. A separate Linux full-RAM experiment with a live packaged Inspector host carved all nine exact-wire-present MCP messages and found no MCP labels in two same-modality non-MCP controls, but missed one MCP classification.

These conditional observations and the constructed generic-JSON-RPC counterexample preclude a general accuracy claim. More host families, repeated controls, field samples, and current-version Windows experiments are needed.

\begin{thebibliography}{99}
\bibitem{mcp-spec} Model Context Protocol, ``Specification, 2026-07-28,'' 2026. [Online]. Available: \url{https://modelcontextprotocol.io/specification/2026-07-28}
\bibitem{jsonrpc} JSON-RPC Working Group, ``JSON-RPC 2.0 Specification.'' [Online]. Available: \url{https://www.jsonrpc.org/specification}
\bibitem{mcprecon-paper} A. Satter, M. Salmon, L. Muhanna, T. T. Spinosa, T. Gharaibeh, and I. Baggili, ``With or Without Logs: Memory Forensic Reconstruction of Model Context Protocol (MCP) Activity in Agentic LLM Systems,'' \emph{Forensic Science International: Digital Investigation}, 2026. [Online]. Available: \url{https://dfrws.org/wp-content/uploads/2026/04/DFRWS_USA_2026_Camera_Ready_Paper_82.pdf}
\bibitem{mcprecon-code} BiT Lab, ``MCPRecon,'' commit 4811d04, 2026. [Online]. Available: \url{https://github.com/BiTLab-BaggiliTruthLab/MCPRecon/tree/4811d04e6c6a522759ade5e2f3fef37d33c6c27e}
\bibitem{volatility} Volatility Foundation, ``Volatility 3 Documentation.'' [Online]. Available: \url{https://volatility3.readthedocs.io/}
\bibitem{lsp} Microsoft, ``Language Server Protocol Specification 3.17.'' [Online]. Available: \url{https://microsoft.github.io/language-server-protocol/specifications/lsp/3.17/specification/}
\bibitem{inspector} Model Context Protocol, ``MCP Inspector mcpdo,'' v2.15.0, 2026. [Online]. Available: \url{https://github.com/modelcontextprotocol/inspector/tree/main/clients/mcpdo}
\bibitem{mcptrace-case} T. Kwon, ``Worked image-to-wire audit: node-d30-r2,'' 2026. [Online]. Available: \url{https://github.com/Kowntaewook/MCPTrace/blob/research/isdfs-mcprecon-comparison/validation/mcprecon_comparison/ground_truth/worked_case_node_d30_r2.md}
\bibitem{mcptrace} T. Kwon, ``MCPTrace,'' source commit ded4c00 and research validation branch, 2026. [Online]. Available: \url{https://github.com/Kowntaewook/MCPTrace/tree/research/isdfs-mcprecon-comparison}
\end{thebibliography}

\end{document}
